\paragrafos{Δυνάμεις με εκθέτη ακέραιο: προς τα αριστερά του μηδενός}{Α.7.9, «Δυνάμεις ρητών αριθμών με εκθέτη ακέραιο»}

\begin{anagke}
Η Πολυξένη γράφει στον πίνακα τις δυνάμεις του $2$: $2^4=16$, $2^3=8$, $2^2=4$, $2^1=2$. Ο Στέλιος κοιτάζει και λέει: «Κάθε φορά που ο εκθέτης μειώνεται κατά $1$, ο αριθμός γίνεται το μισό: $16$, $8$, $4$, $2$». «Και αν συνεχίσουμε;» ρωτάει η Αγγελική. Ο επόμενος αριθμός πρέπει να είναι το μισό του $2$, δηλαδή $1$: $2^0=1$. Και μετά το μισό του $1$: $\frac{1}{2}$, δηλαδή $2^{-1}=\frac{1}{2}$. Και $2^{-2}=\frac{1}{4}$, $2^{-3}=\frac{1}{8}$.

«Αρνητικός εκθέτης;» αναρωτιέται ο Βασίλης. «Δεν μπορώ να πολλαπλασιάσω τον $2$ μείον τρεις φορές!». Η κυρία Αλεξάνδρα εξηγεί: «Δεν χρειάζεται να το κάνετε. Αφού ο εκθέτης που μικραίνει διαιρεί με το $2$, μπορούμε να ορίσουμε ότι συνεχίζεται το ίδιο μοτίβο και στους αρνητικούς, ώστε να ισχύουν οι ιδιότητες των δυνάμεων που ξέρετε». Ο Νίκος το συνδέει με τις μονάδες: «Ένα χιλιοστό είναι το $\frac{1}{1\,000}$ του μέτρου, δηλαδή $10^{-3}$ μέτρα. Αυτό θα γράφουμε στο εργαστήριο».

Στην παράγραφο αυτή η παρέα αποκτά ένα εργαλείο για τους πολύ μεγάλους και τους πολύ μικρούς αριθμούς, και δίνει νόημα στο $\alpha^0$ και στο $\alpha^{-\nu}$.
\end{anagke}

\begin{lexi}{αρνητικός εκθέτης}{αντίστροφος, δύναμη}
Στα μαθηματικά οι νέες έννοιες ορίζονται ώστε να «συνεχίζουν το μοτίβο» και να ισχύουν οι παλιοί κανόνες. Η δύναμη με αρνητικό εκθέτη δεν σημαίνει «πολλαπλασιάζω αρνητικές φορές» αλλά «αντίστροφος της δύναμης»: το $2^{-3}$ είναι ο \lexiword{αντίστροφος} του $2^3$.
\end{lexi}

\begin{metafrasi}
Το μισό του $2$. & $2^0=1$\\
Το μισό του $1$. & $2^{-1}=\dfrac{1}{2}$\\
Ένα χιλιοστό του μέτρου. & $10^{-3}$ μ.\\
Η δύναμη με αρνητικό εκθέτη. & $\alpha^{-\nu}=\dfrac{1}{\alpha^{\nu}}$\\
Κλάσμα με αρνητικό εκθέτη. & $\left(\dfrac{\alpha}{\beta}\right)^{-\nu}=\left(\dfrac{\beta}{\alpha}\right)^{\nu}$\\
\end{metafrasi}

\tmima{Η θεωρία}

Ας γράψουμε τις δυνάμεις του $2$ και ας συνεχίσουμε προς τα δεξιά:
\begin{center}
\renewcommand{\arraystretch}{1.3}
\begin{tabular}{c|ccccccc}
$\nu$ & $3$ & $2$ & $1$ & $0$ & $-1$ & $-2$ & $-3$\\\hline
$2^\nu$ & $8$ & $4$ & $2$ & $1$ & $\frac{1}{2}$ & $\frac{1}{4}$ & $\frac{1}{8}$
\end{tabular}
\end{center}
Κάθε φορά που ο εκθέτης μειώνεται κατά $1$, η τιμή διαιρείται με $2$. Έτσι ορίζουμε:

\begin{orismos}{δύναμη με εκθέτη ακέραιο}
Για κάθε ρητό $\alpha\neq0$ και φυσικό $\nu$:
$$\alpha^0=1,\qquad \alpha^{-\nu}=\dfrac{1}{\alpha^{\nu}}.$$
\end{orismos}

Για παράδειγμα $5^0=1$, $(-7)^0=1$, $3^{-2}=\frac{1}{9}$, $(-2)^{-3}=\frac{1}{(-2)^3}=-\frac{1}{8}$, $10^{-3}=\frac{1}{1\,000}=0{,}001$. Το $0^0$ και τα $0^{-\nu}$ δεν ορίζονται.

\tmimasmall{Κλάσμα στη δύναμη}

Αν η βάση είναι κλάσμα, ο αρνητικός εκθέτης αναποδογυρίζει το κλάσμα:
$$\left(\dfrac{\alpha}{\beta}\right)^{-\nu}=\left(\dfrac{\beta}{\alpha}\right)^{\nu}\quad(\alpha,\beta\neq0).$$
Για παράδειγμα $\left(\frac{1}{2}\right)^{-3}=2^3=8$ και $\left(\frac{2}{3}\right)^{-2}=\left(\frac{3}{2}\right)^2=\frac{9}{4}$.

\begin{idiotita}{ιδιότητες με ακέραιο εκθέτη}
Οι ιδιότητες των δυνάμεων ισχύουν και για ακέραιους εκθέτες ($\alpha,\beta\neq0$): $\alpha^\mu\cdot\alpha^\nu=\alpha^{\mu+\nu}$, $\alpha^\mu:\alpha^\nu=\alpha^{\mu-\nu}$, $(\alpha\beta)^\nu=\alpha^\nu\beta^\nu$, $\left(\frac{\alpha}{\beta}\right)^\nu=\frac{\alpha^\nu}{\beta^\nu}$, $(\alpha^\mu)^\nu=\alpha^{\mu\nu}$.
\end{idiotita}

Με το νέο ορισμό η διαίρεση δεν χρειάζεται περιορισμό $\mu>\nu$: $5^3:5^5=5^{-2}=\frac{1}{25}$. Και βγαίνει τώρα φυσιολογικά ότι $\alpha^0=1$: $\alpha^3:\alpha^3=\alpha^{3-3}=\alpha^0$ και, αφού $\alpha^3:\alpha^3=1$, είναι $\alpha^0=1$.

\begin{prosoxi}
\begin{itemize}[leftmargin=1.2em,itemsep=0.3ex,topsep=2pt]
\item Το $2^{-3}$ δεν είναι $-8$ ούτε $-6$: είναι $\frac{1}{8}$, θετικό. Το αρνητικό σημείο στον εκθέτη δεν αλλάζει το πρόσημο.
\item $\alpha^0=1$ και όχι $0$ (για $\alpha\neq0$).
\item $\left(\frac{1}{2}\right)^{-3}=8$, όχι $\frac{1}{8}$: αναποδογυρίζω το κλάσμα και κάνω τον εκθέτη θετικό.
\item Το πρόσημο της βάσης παίζει ρόλο: $(-2)^{-3}=-\frac{1}{8}$ και $(-2)^{-2}=\frac{1}{4}$.
\item Το $0^{-1}$ δεν ορίζεται (θα σήμαινε $\frac{1}{0}$).
\end{itemize}
\end{prosoxi}

\tmima{Λυμένα παραδείγματα}

\paradeigma{Συνεχίζω τον πίνακα}
\tixero Οι δυνάμεις του $3$: $3^3=27$, $3^2=9$, $3^1=3$.\par
\tithelo Να βρω τα $3^0$, $3^{-1}$, $3^{-2}$, $3^{-3}$.\par
\possindeo Κάθε φορά ο εκθέτης μειώνεται και η τιμή διαιρείται με $3$: $3:3=1$, $1:3=\frac{1}{3}$, $\frac{1}{3}:3=\frac{1}{9}$, $\frac{1}{9}:3=\frac{1}{27}$.\par
\apantisi $3^0=1$, $3^{-1}=\frac{1}{3}$, $3^{-2}=\frac{1}{9}$, $3^{-3}=\frac{1}{27}$.\par
\elegxos Σύμφωνα με τον ορισμό $3^{-2}=\frac{1}{3^2}=\frac{1}{9}$.

\paradeigma{Υπολογισμοί}
\tixero Οι $3^{-2}$, $(-2)^{-3}$, $\left(\frac{1}{2}\right)^{-3}$, $\left(\frac{2}{3}\right)^{-2}$, $5^0$.\par
\tithelo Να τους υπολογίσω.\par
\possindeo $3^{-2}=\frac{1}{9}$. $(-2)^{-3}=\frac{1}{(-2)^3}=-\frac{1}{8}$. $\left(\frac{1}{2}\right)^{-3}=2^3=8$. $\left(\frac{2}{3}\right)^{-2}=\left(\frac{3}{2}\right)^2=\frac{9}{4}$. $5^0=1$.\par
\apantisi $\frac{1}{9}$, $-\frac{1}{8}$, $8$, $\frac{9}{4}$, $1$.\par
\elegxos $\left(\frac{1}{2}\right)^{-3}\cdot\left(\frac{1}{2}\right)^{3}=\left(\frac{1}{2}\right)^0=1$ και $8\cdot\frac{1}{8}=1$.

\paradeigma{Με τις ιδιότητες}
\tixero Οι $2^{-3}\cdot2^5$, $10^4:10^{-2}$, $(3^{-1})^2$, $(2\cdot5)^{-2}$.\par
\tithelo Να τους υπολογίσω.\par
\possindeo $2^{-3+5}=2^2=4$. $10^{4-(-2)}=10^6$. $3^{-2}=\frac{1}{9}$. $2^{-2}\cdot5^{-2}=\frac{1}{4}\cdot\frac{1}{25}=\frac{1}{100}$.\par
\apantisi $4$, $10^6=1\,000\,000$, $\frac{1}{9}$, $\frac{1}{100}$.\par
\elegxos $(2\cdot5)^{-2}=10^{-2}=\frac{1}{100}$.

\paradeigma{Παράσταση}
\tixero Η παράσταση $\left(\frac{1}{2}\right)^{-2}+3^0-2^{-1}$.\par
\tithelo Να την υπολογίσω.\par
\possindeo $\left(\frac{1}{2}\right)^{-2}=2^2=4$, $3^0=1$, $2^{-1}=\frac{1}{2}$. Άρα $4+1-\frac{1}{2}=\frac{9}{2}$.\par
\apantisi $\frac{9}{2}$.\par
\elegxos $4{,}5=4+1-0{,}5$.

\paradeigma{Εξίσωση με δύναμη}
\tixero Η εξίσωση $2^x=\frac{1}{16}$.\par
\tithelo Να βρω το $x$.\par
\possindeo Γράφω $\frac{1}{16}=\frac{1}{2^4}=2^{-4}$. Άρα $x=-4$.\par
\apantisi $x=-4$.\par
\elegxos $2^{-4}=\frac{1}{16}$.

\tmima{Ασκήσεις}

\zoni{Το κατάλαβα}{αρνητικοί εκθέτες και μηδέν}

\ex Υπολόγισε: α) $2^{-1}$ \quad β) $2^{-3}$ \quad γ) $10^{-2}$ \quad δ) $5^{-2}$ \quad ε) $4^0$ \quad στ) $(-3)^0$.

\ex Γράψε ως δεκαδικούς: α) $10^{-1}$ \quad β) $10^{-3}$ \quad γ) $10^{-5}$.

\ex Υπολόγισε: α) $(-2)^{-3}$ \quad β) $(-2)^{-2}$ \quad γ) $(-1)^{-5}$ \quad δ) $\left(\frac{1}{2}\right)^{-3}$ \quad ε) $\left(\frac{2}{3}\right)^{-2}$ \quad στ) $\left(-\frac{3}{4}\right)^{-1}$.

\ex Γράψε ως δύναμη με αρνητικό εκθέτη: α) $\frac{1}{9}$ (με βάση $3$) \quad β) $\frac{1}{32}$ (βάση $2$) \quad γ) $0{,}0001$ (βάση $10$) \quad δ) $\frac{1}{(-5)^3}$.

\ex Υπολόγισε με τις ιδιότητες: α) $2^{-3}\cdot2^5$ \quad β) $10^4:10^{-2}$ \quad γ) $(3^{-1})^2$ \quad δ) $(2\cdot5)^{-2}$ \quad ε) $7^{-4}\cdot7^4$ \quad στ) $(-2)^{-3}\cdot(-2)^5$.

\ex Υπολόγισε: α) $\left(\frac{1}{2}\right)^{-2}+3^0-2^{-1}$ \quad β) $2^{-2}+2^{-1}$ \quad γ) $10^2\cdot10^{-5}\cdot10^3$ \quad δ) $(3^2)^{-1}$ \quad ε) $5^0\cdot(-2)^{-2}$.

\ex Βρες το $x$: α) $2^x=\frac{1}{16}$ \quad β) $3^x=\frac{1}{27}$ \quad γ) $10^x=0{,}001$ \quad δ) $\left(\frac{1}{2}\right)^x=8$ \quad ε) $5^x=1$.

\ex Σωστό ή λάθος; α) $2^{-1}=-2$. \quad β) $\alpha^0=1$ για $\alpha\neq0$. \quad γ) $\left(\frac{1}{2}\right)^{-1}=2$. \quad δ) $10^{-3}=0{,}003$. \quad ε) $\alpha^{-\nu}=\frac{1}{\alpha^\nu}$. \quad στ) $2^{-2}\cdot2^2=1$. \quad ζ) $(-2)^{-2}=-\frac{1}{4}$.

\zoni{Το χρησιμοποιώ}{μονάδες, μισά, αξίες}

\ex \textbf{Οι μονάδες.} Γράψε με δυνάμεις του $10$ σε μέτρα: $1$ χιλιόμετρο, $1$ εκατοστό, $1$ χιλιοστό, $1$ μικρόμετρο ($0{,}000001$ μ.), $1$ νανόμετρο ($0{,}000000001$ μ.).

\ex \textbf{Η πίτα.} Κόβουμε κάθε φορά το μισό από ό,τι έμεινε. Τι μέρος της πίτας έχει μείνει μετά από $5$ κοψίματα; Γράψε το ως δύναμη του $2$.

\ex \textbf{Το ραδιενεργό υλικό.} Ένα υλικό 64 g χάνει το μισό της μάζας του κάθε ημέρα. Πόσο θα είναι μετά από $3$ ημέρες; Πόσο ήταν πριν από $2$ ημέρες;

\ex \textbf{Ο πίνακας.} Συμπλήρωσε τις δυνάμεις του $3$ από $3^4$ έως $3^{-3}$.

\ex \textbf{Το κινητό.} Η αξία ενός κινητού των $400$ € μειώνεται στο μισό κάθε χρόνο. Πόση αξία έχει μετά από $3$ έτη; Γράψε το με δύναμη.

\zoni{Το σκέφτομαι}{γιατί το μηδέν δίνει ένα}

\ex Γιατί $\alpha^0=1$ και όχι $0$; Χρησιμοποίησε το μοτίβο των δυνάμεων του $2$.

\ex Δείξε ότι $\left(\frac{2}{3}\right)^{-2}=\left(\frac{3}{2}\right)^2$.

\ex Σε τι διαφέρουν τα $2^{-3}$ και $-2^3$;

\ex Αν ισχύει $\alpha^\mu\cdot\alpha^\nu=\alpha^{\mu+\nu}$ και $\mu=-3$, $\nu=3$, τι σημαίνει $\alpha^0$;

\ex Γιατί δεν ορίζεται το $0^{-1}$;

\ex Υπολόγισε τα $(-1)^{-1}$, $(-1)^{-2}$, $(-1)^{-3}$ και βρες τον κανόνα.

\zoni{Για το σπίτι}{πόσο λεπτό είναι ένα φύλλο;}

\ex Ένα φύλλο χαρτιού έχει πάχος περίπου $0{,}1$ mm. Γράψε το πάχος σε μέτρα με δύναμη του $10$.
